\documentclass[11pt,a4paper]{article}
\usepackage[margin=25mm]{geometry}
\usepackage[T1]{fontenc}
\usepackage{lmodern,microtype}
\usepackage{amsmath,amssymb,amsthm,bm}
\usepackage{booktabs,longtable,tabularx,array}
\usepackage{xurl,hyperref}
\usepackage{listings}
\hypersetup{colorlinks=true,linkcolor=blue,urlcolor=blue}
\lstset{basicstyle=\ttfamily\small,breaklines=true,columns=fullflexible}
\setlength{\emergencystretch}{3em}
\newcommand{\E}{\mathbf E}
\newcommand{\HH}{\mathbf H}
\newcommand{\n}{\mathbf n}
\newcommand{\curl}{\nabla\times}
\newcommand{\divg}{\nabla\cdot}
\newcommand{\R}{\mathbb R}
\newcommand{\ip}[2]{\left(#1,#2\right)}
\newcommand{\src}[3]{\par\smallskip\noindent\textbf{Source.} \path{#1}, \path{#2}, lines #3.\par\smallskip}
\newcommand{\kernel}{\texttt{src/solver/kernels/}}
\newtheorem{proposition}{Proposition}[section]
\title{The Implemented Three-Dimensional Maxwell Poisson-Bracket Operator in DiscoGMPI\\[0.5em]\large A source-based reconstruction and independent algebraic audit}
\author{Implementation snapshot: branch \texttt{main} with local changes}
\date{21 September 2026}
\begin{document}
\maketitle
\begin{abstract}
This report reconstructs the spatial Maxwell operator actually executed by
DiscoGMPI on affine tetrahedra, rather than prescribing a generic discontinuous
Galerkin method. Starting with the monomial Gram construction, nodal bases,
physical mappings and face restrictions, it derives the asymmetric electric
and magnetic volume terms, centered and orientation-selected alternating
traces, mass-weighted global system, and finite-dimensional bracket. It then
identifies the assumptions under which this bracket is constant and skew,
and separates those mathematical properties from floating-point identities,
time integration, and monitored divergence diagnostics. MPI realization and
boundary forcing are traced to executable routines. In particular, different
dispatch paths do not always honor the same flux request: a material-aware
wrapper defaults to centered traces even when alternating flux is requested.
Independent small-matrix and distributed checks accompany the reconstruction;
no production operator has been altered.
\end{abstract}
\tableofcontents
\newpage

\section{Scope, evidence and executable path}\label{sec:scope}
The baseline commit is
\path{e5bd7abe249efab7806acfb197af1141313523f9}, on branch \texttt{main}.
The working tree is not clean. In particular the nodal/face constructors,
distributed DG support, diagnostics, driver support and example drivers have
local changes; magnetic initialization is an untracked source addition at
inspection time. Thus the commit alone cannot reproduce this report.
The accompanying provenance file records source hashes and working-tree state.
No applicable \path{AGENTS.md} was found in the repository or its ancestor
directories. Line references refer to this working-tree snapshot.

Three evidence classes are used throughout. \emph{Source} identifies executed
array operations and dispatch. \emph{Derivation} identifies consequences in
exact real arithmetic under explicitly stated assumptions. \emph{Verification}
identifies finite, executed numerical checks, whose domains do not constitute
general proofs. Observed inconsistencies are retained, not silently corrected.

The main thread is \path{PoissonBracketFormulation}, not
\path{HesthavenWarburtonFormulation}. Both coexist in the same module. The
latter uses strong volume curls for both fields and permits upwind penalties;
the former uses an adjoint magnetic volume operator and permits centered or
alternating interface fluxes on its homogeneous-coefficient path. There is
no active \path{rhs_hamiltonian_H!} or \path{FaceCommPlan} API on this path.
\path{src/TraceMaps.jl} contains a separate generic trace-exchange facility;
the Maxwell runtime instead uses \path{DistributedMaxwellExchange}.
\src{src/DiscoGMPI.jl}{include chain}{352--417}
\src{src/solver/data_containers.jl}{AbstractMaxwellDGFormulation; PoissonBracketFormulation}{438--454}

\subsection{Representative call graph}
The periodic driver provides a particularly direct, unforced realization:
\begin{lstlisting}
examples/distributed_periodic_poisson_bracket_maxwell.jl
  main -> run_experiment -> load_or_prepare_distributed_dg
    load_periodic_box_mesh -> read_mesh_file_tet_vtk
    prepare_distributed_mesh_partition[_collective]
    build_distributed_dg_from_partition
      build_distributed_dg_from_local_mesh -> DGDiscretization
        topology -> geometry -> mappings -> reference tet
        face operators -> physical operators -> trace maps -> flux faces
  build_distributed_periodic_maxwell_exchange
  interpolate_maxwell_field [optional curl-potential replacement of H]
  distributed_periodic_partitioned_symplectic_rk_step!
    partitioned_symplectic_rk_step! -> RHS callback
      maxwell_rhs_periodic!
        ordinary ghost exchange; periodic trace exchange
        volume; ordinary faces; periodic faces; physical boundaries
        zero_ghost_maxwell_rhs!
\end{lstlisting}
The representative driver calls the scalar \(\varepsilon,\mu\) overload, not
the element-material overload. Its optional curl initializer does not change
that evolution dispatch. Its diagnostic cubature setting does not control
the evolution mass or stiffness matrices.
\src{examples/distributed_periodic_poisson_bracket_maxwell.jl}{load_or_prepare_distributed_dg; run_experiment}{534--595, 704--753, 1006--1018}
\src{src/solver/kernels/dg_discretization.jl}{DGDiscretization}{5--42}

\section{Continuous fields and conventions}\label{sec:continuous}
Consider a stationary mesh of affine tetrahedra \(\mathcal T_h\) covering
\(\Omega\subset\R^3\). Materials are positive, time-independent scalar
constants on each element, \(\varepsilon_K,\mu_K>0\). The baseline source-free
equations represented by the conservative spatial operator are
\begin{equation}\label{eq:maxwell}
 \varepsilon\partial_t\E=\curl\HH,\qquad
 \mu\partial_t\HH=-\curl\E,\qquad
 \mathbf D=\varepsilon\E,\quad\mathbf B=\mu\HH.
\end{equation}
There is no general volume current supplied to the baseline RHS. The PML adds
separate source terms, and incident PEC data add boundary forcing; neither is
part of the closed homogeneous Hamiltonian system below.

\begin{longtable}{p{0.24\linewidth}p{0.69\linewidth}}
\toprule Symbol & Meaning\\\midrule
\(\widehat K,\widehat x=(r,s,t)\) & Reference tetrahedron and reference coordinates; physical time is denoted \(\tau\) when confusion with coordinate \(t\) is possible.\\
\(K,x=\Phi_K(\widehat x)\) & Physical element and its affine map.\\
\(\n_K\), \(\n\) & Outward physical unit normal of an element; on a two-sided face, \(\n\) points out of the chosen minus element.\\
\(v^-,v^+\) & Values at the same physical face point, or translated periodic point, in minus ordering.\\
\(\Delta v=v^--v^+\) & Componentwise scalar/vector difference; no hidden normal factor.\\
\(j_n(v)=\n\cdot\Delta v\) & Normal jump relative to the chosen normal.\\
\(j_t(v)=\n\times\Delta v\) & Tangential vector jump, with this cross-product order.\\
\(\{v\}=(v^-+v^+)/2\) & Arithmetic mean, not an impedance-weighted mean.\\
\(\ip{u}{v}_K\) & \(\int_K u\cdot v\,dx\), or ordinary scalar product for scalar functions.\\
\(\langle u,v\rangle_f\) & Face integral \(\int_f u\cdot v\,dS\).\\
\(R_f,P_f\) & Nodal face restriction and plus-to-minus permutation matrix.\\
\(N(\n)\) & The \(3\times3\) matrix such that \(N(\n)v=\n\times v\); \(N^T=-N\).\\
\(M_0,M_E,M_H\) & Global unweighted vector mass and its electric/magnetic material-weighted counterparts.\\
\bottomrule
\end{longtable}
The curl sign follows
\((\partial_y v_z-\partial_z v_y,\partial_z v_x-\partial_x v_z,
\partial_x v_y-\partial_y v_x)\). In particular,
\begin{equation}\label{eq:curlgreen}
 (\curl u,v)_K=(u,\curl v)_K+\langle\n_K\times u,v\rangle_{\partial K}.
\end{equation}
This fixes the signs of every subsequent boundary term.

\subsection{Hamiltonian and functional derivatives}
The electromagnetic Hamiltonian is
\begin{equation}\label{eq:contenergy}
 \mathcal H=\frac12\int_\Omega(\varepsilon|\E|^2+\mu|\HH|^2)\,dx.
\end{equation}
For a differentiable functional \(F\), let \(F_E,F_H\) be its
\emph{unweighted} \(L^2\) gradients. Its material-weighted gradients are
\(F_E^w=\varepsilon^{-1}F_E\), \(F_H^w=\mu^{-1}F_H\).
Thus \(\mathcal H_E=\varepsilon\E\),
\(\mathcal H_H=\mu\HH\), while \(\mathcal H_E^w=\E\),
\(\mathcal H_H^w=\HH\). A relevant continuous bulk bracket is
\begin{equation}\label{eq:contbracket}
 \{F,G\}_{\rm bulk}=
 (F_E^w,\curl G_H^w)_\Omega
 -(F_H^w,\curl G_E^w)_\Omega.
\end{equation}
The symmetric part is a boundary pairing:
\begin{equation}\label{eq:bracketboundary}
 \{F,G\}_{\rm bulk}+\{G,F\}_{\rm bulk}
 =\int_{\partial\Omega}\n\cdot
 (G_H^w\times F_E^w+F_H^w\times G_E^w)\,dS.
\end{equation}
Periodicity or compatible homogeneous PEC/PMC domains for the gradients
remove this term. Otherwise one needs boundary power/ports, not a claim of
a closed antisymmetric bulk bracket. For smooth solutions,
\begin{equation}
 \dot{\mathcal H}=-\int_{\partial\Omega}\n\cdot(\E\times\HH)\,dS.
\end{equation}
For discontinuous materials these equations and integrations are interpreted
piecewise, with compatible interface transmission. The DG bracket will be
defined from its actual matrix blocks, not inferred merely from
\eqref{eq:contbracket}.

\section{Reference polynomials and nodal representation}\label{sec:reference}
The reference vertices, in \emph{geometric vertex numbering}, are
\begin{equation}
 v_1=(-1,-1,-1),\quad v_2=(1,-1,-1),\quad
 v_3=(-1,1,-1),\quad v_4=(-1,-1,1).
\end{equation}
The volume is \(4/3\). Face numbering and oriented vertex triples are
\begin{center}
\begin{tabular}{ccccc}\toprule
Face & Vertices & Plane & Reference outward normal & Area\\\midrule
1 & (2,3,4) & \(r+s+t=-1\) & \((1,1,1)/\sqrt3\) & \(2\sqrt3\)\\
2 & (1,4,3) & \(r=-1\) & \((-1,0,0)\) & 2\\
3 & (1,2,4) & \(s=-1\) & \((0,-1,0)\) & 2\\
4 & (1,3,2) & \(t=-1\) & \((0,0,-1)\) & 2\\\bottomrule
\end{tabular}
\end{center}
\src{src/solver/data_containers.jl}{TET_FACES; REF_TET_VERTEX_COORDS; REF_TET_FACE_VERTEX_IDS}{27--32, 152--164}

The scalar space is the complete total-degree space \(P_p(\widehat K)\),
not a tensor-product polynomial space. Its volume and face dimensions are
\begin{equation}\label{eq:dimensions}
 N_p=\binom{p+3}{3},\qquad N_{fp}=\binom{p+2}{2}.
\end{equation}
Degree zero is supported by the reference-volume constructor, but the
ordinary face-node-extraction discretization requires nodes on faces. Its
centroid-only degree-zero node cannot supply those traces; this report does
not claim a working production \(p=0\) DG face discretization.

\subsection{Equispaced and warp-blend nodes}
The default equispaced loop is \(i=0{:}p\), then \(j=0{:}p-i\), then
\(k=0{:}p-i-j\), with \(l=p-i-j-k\). It sets
\begin{equation}
 (\lambda_1,\lambda_2,\lambda_3,\lambda_4)=(i,j,k,l)/p,
 \qquad(r,s,t)=(-1+2\lambda_2,-1+2\lambda_3,-1+2\lambda_4).
\end{equation}
Consequently the \(p=1\) solver nodes are \(v_4,v_3,v_2,v_1\), not
\(v_1,v_2,v_3,v_4\). For \(p=0\) the node is \((-1/2,-1/2,-1/2)\).
\src{src/solver/kernels/reference_tet.jl}{num_tet_nodes; equispaced_tet_nodes}{6--55}

The alternative \path{warp_blend_tet_nodes} uses its own equispaced seed
ordering, transforms to an equilateral tetrahedron, accumulates face shifts
in two tangential directions, and maps back. Edge targets are
Gauss--Lobatto--Legendre points: endpoints and roots of \(P_p'\), computed
by Newton iteration. The degree-dependent blending parameter is the stored
\path{WARP_BLEND_TET_ALPHA} table (degrees 1--15), with fallback 1 above
that range. The construction does not sort its output into the default
equispaced ordering.

More explicitly, the two-dimensional face shift combines the three edge
warps with factors \(4L_bL_c(1+(\alpha L_a)^2)\) in directions at angles
\(0,2\pi/3,4\pi/3\). For a tetrahedral face opposite barycentric coordinate
\(L_a\), its interior extension uses
\begin{equation}
 b_f=\frac{(1+(\alpha L_a)^2)L_bL_cL_d}
 {(L_b+L_a/2)(L_c+L_a/2)(L_d+L_a/2)}.
\end{equation}
The code treats small denominators and face edges separately with tolerance
\(10^{-10}\); the implemented branch logic, not this rational expression
alone, defines edge values. These are interpolation nodes, not volume
integration points.
\src{src/solver/kernels/reference_tet.jl}{gauss_lobatto_legendre_nodes; eval_shift_2d; warp_blend_tet_nodes}{115--230, 276--388}

\subsection{The actual modal basis: Gram orthonormalization}
The internal basis is not implemented as a direct Jacobi/Dubiner recurrence.
Let \(m_a=r^{a_1}s^{a_2}t^{a_3}\), ordered first by total degree, then
increasing \(a_1\), then increasing \(a_2\). The exact-in-real-arithmetic
monomial Gram matrix is \(G_{ab}=\int_{\widehat K}m_am_b\).
For unit-simplex coordinates \((u,v,w)\),
\begin{equation}
 \int_{u,v,w\ge0,\ u+v+w\le1}u^iv^jw^k\,du\,dv\,dw
 =\frac{i!j!k!}{(i+j+k+3)!}.
\end{equation}
Binomial expansion of \(r=2u-1,s=2v-1,t=2w-1\), multiplied by 8,
gives the moments used in \(G\). Computation is Float64 even though factorials
are first formed as big integers. Cancellation in these moments can therefore
limit high-degree conditioning.

The Cholesky factor is \(G=U^TU\). With \(C=U^{-1}\), the modal row vector is
\(\psi(\widehat x)=m(\widehat x)C\), so
\begin{equation}
 C^TGC=I,\qquad V_{ia}=\psi_a(\widehat x_i),\qquad
 \ell(\widehat x)=\psi(\widehat x)V^{-1}.
\end{equation}
Hence \(\ell_j(\widehat x_i)=\delta_{ij}\), nodal values \(u\) map to modal
coefficients \(V^{-1}u\), and quadrature values are \(V_qV^{-1}u\).
The matrices explicitly store both \(V\) and \(V^{-1}\).
\src{src/solver/kernels/reference_tet.jl}{monomial_exponents_3d; ref_tet_monomial_integral; build_orthonormal_tet_basis}{426--440, 522--582, 616--646}

Differentiating the monomials before transforming gives
\begin{equation}\label{eq:refmatrices}
 D_r=V_rV^{-1},\quad D_s=V_sV^{-1},\quad D_t=V_tV^{-1},
 \quad M=V^{-T}V^{-1},\quad S_a=MD_a.
\end{equation}
Thus \((S_a)_{ij}=\int\ell_i\partial_a\ell_j\): row \(i\) is the test index,
column \(j\) the trial index. The derivative is on the \emph{column} basis.
All these matrices are \(N_p\times N_p\). The older nested-loop analytic
mass/stiffness helpers remain in the file, but are not invoked by
\path{build_reference_tet}.
\src{src/solver/kernels/reference_tet.jl}{orthonormal_grad_vandermonde_tet; build_reference_tet}{596--614, 769--848}

\section{Face integration, restriction and lifting}\label{sec:faces}
For a face \(f\), \(R_f\in\R^{N_{fp}\times N_p}\) selects the ordered
face nodes. On equispaced meshes they are reordered to the reference
triangle's order. On warp-blend faces, a separate triangular basis is built
from the actual extracted face coordinates. The mappings to the reference
triangle of area 2 are \((s,t),(t,s),(r,t),(s,r)\), respectively for faces
1 through 4. This is why a face-node permutation must be applied consistently
to values and face mass matrices.

The triangular modal Gram/Cholesky construction gives
\(M_\triangle=V_\triangle^{-T}V_\triangle^{-1}\). The active reference face
matrix is
\begin{equation}\label{eq:facemass}
 W_f=\frac{|\widehat f|}{2}M_{\triangle,f},\quad
 B_f=R_f^TW_fR_f,\quad L_f=M^{-1}R_f^TW_f.
\end{equation}
Here \(W_f\) is \(N_{fp}\times N_{fp}\), \(B_f\) is \(N_p\times N_p\), and
the precomputed \path{face_lift[f]} is \(N_p\times N_{fp}\).
The additional \path{Emat} and \path{LIFT} fields hold
\(\sum_fB_f\) and \(M^{-1}\sum_fB_f\). Maxwell face loops use the individual
\(L_f\), not the combined matrix.
\src{src/solver/kernels/face_operators.jl}{build_reference_tri_from_nodes; build_reference_face_mass_from_tri}{644--684, 734--737}
\src{src/solver/kernels/face_operators.jl}{ordered_reference_face_nodes; build_reference_face_operators}{792--848, 871--972}
\src{src/solver/kernels/face_operators.jl}{build_reference_face_lift_matrices}{243--268}

\subsection{Which quadrature is actually used?}
The baseline volume and active face matrices above use analytic polynomial
moments through orthonormalization; they are not nodal mass lumping or a
Jaskowiec--Sukumar assembly. The mass is full, symmetric positive definite
in exact arithmetic, and is factored/solved, not approximated as diagonal.
The volume RHS factors \(M\) with Cholesky once per RHS invocation, before
the element loop; that factor is not cached across calls on this path.
Face lifting solves against \(M\) during construction and subsequently
multiplies precomputed matrices, with no mass solve inside the face loop.

A separate face quadrature implementation uses a Duffy transform of a tensor
Gauss--Legendre rule, with \(n=\lceil(d+2)/2\rceil\) points per direction
for target degree \(d\). On \([0,1]^2\), \((u,v)\mapsto(u,(1-u)v)\), so
weights acquire \(1-u\). Their sum is \(1/2\), then multiplying by
\(2|\widehat f|\) gives surface weights. Taking \(d=2p\) is polynomially
exact for face basis products. However, its old nodal-evaluation helper
currently throws a shape error in the audit; it is not the production
face constructor. The audit independently evaluates the modal matrix at the
same quadrature points to check \eqref{eq:facemass}.
\src{src/solver/kernels/face_operators.jl}{triangle_quadrature_gauss; reference_face_quadrature; eval_nodal_basis_at_points}{80--137, 173--240}

Diagnostic volume cubature uses tabulated Jaskowiec--Sukumar rules of orders
2--20. Tabulated coordinates are transformed by \(2x-1\), and weights by
\(8/6\); final weights sum to \(4/3\). Physical diagnostic weights are
\(|\det J_K|w_q\). The default diagnostic order is \(\max(2,2p+4)\), subject
to availability. These rules integrate polynomial products to their stated
degree, not trigonometric analytical errors exactly. Solution nodes and
cubature points are distinct.
\src{src/solver/kernels/JaskowiecSukumar.jl}{get_JaskowiecSukumar_cubature}{1--2694; coordinate/weight transformation at 2683--2694}
\src{src/MaxwellDriverSupport.jl}{resolved_maxwell_cubature_order}{414--430}

\subsection{Discrete integration by parts}
For representable polynomials and exact moment/face integration,
\begin{equation}\label{eq:sbp}
 S_a+S_a^T=\sum_f\widehat n_{f,a}B_f,\qquad a=r,s,t.
\end{equation}
This follows entry by entry from
\(\partial_a(\ell_i\ell_j)\) and the divergence theorem. It is not a
pointwise product rule for arbitrary interpolated nonlinear functions.
The stored Float64 matrices satisfy it only to construction roundoff;
Section~\ref{sec:verification} measures that residual.

\section{Affine geometry and physical fields}\label{sec:geometry}
Let the physical vertices be \(x_1,\ldots,x_4\). The executable map is
\begin{equation}\label{eq:map}
 \Phi_K(r,s,t)=\sum_{a=1}^4\lambda_a(r,s,t)x_a,
 \quad J_K=\frac12[x_2-x_1\;x_3-x_1\;x_4-x_1],
 \quad j_K=|\det J_K|.
\end{equation}
The stored inverse is \(J_K^{-1}\); gradients use \(J_K^{-T}\). Thus
\begin{equation}\label{eq:physicalD}
 D_{K,x}=\sum_{a=1}^3(J_K^{-1})_{a1}D_a,
 \quad D_{K,y}=\sum_a(J_K^{-1})_{a2}D_a,
 \quad D_{K,z}=\sum_a(J_K^{-1})_{a3}D_a.
\end{equation}
Using rows instead of columns of \(J^{-1}\) here would be incorrect.
Signed determinants are retained for orientation diagnostics; positive
\(j_K\) is used for integration. A negative determinant is not, by itself,
rejected: the normal is oriented independently with the owner centroid.
The mapping rejects \(|\det J_K|\le\epsilon_{\rm mach}\), an absolute
degeneracy threshold rather than a scale-invariant quality test.
\src{src/solver/kernels/metrics.jl}{ref_tet_shape_functions; build_tet_mapping; reference_gradient_to_physical}{7--98}

The physical scalar mass is \(M_K=j_KM\). The stored physical stiffness is
\(S_{K,a}=MD_{K,a}\), \emph{without} the determinant; its actual physical
integral is \(j_KS_{K,a}\). The omitted factor cancels with \(M_K\) during
volume inversion. The componentwise vector mass is \(I_3\otimes M_K\).
\src{src/solver/kernels/physical_operators.jl}{physical_derivative_matrices; physical_stiffness_matrices; build_physical_element_operator}{6--63}

For a face, code computes the triangle cross product, area and outward unit
normal, flipping the normal if it points toward the owner centroid. Define
\(s_{Kf}=|f|/|\widehat f|\). Then
\begin{equation}\label{eq:geomlift}
 M_K^{-1}R_f^T(s_{Kf}W_f)g
 =\frac{s_{Kf}}{j_K}L_fg.
\end{equation}
Normals are \emph{unit}, not area-scaled; area appears only in
\(s_{Kf}\). The signed map and outward orientation imply the physical version
\(j_K(S_{K,a}+S_{K,a}^T)=\sum_fs_{Kf}n_{Kf,a}B_f\).
\src{src/solver/kernels/geometry.jl}{raw_triangle_normal_area; oriented_face_geometry}{113--148}
\src{src/solver/kernels/maxwell.jl}{add_lifted_face_contribution!}{1171--1220}

The vector space is \(V_h=[P_p(\mathcal T_h)]^3\), with Cartesian components
mapped as scalar polynomials:
\begin{equation}
 \E_h|_K=\sum_{i=1}^{N_p}\mathbf e_{K,i}\ell_i\circ\Phi_K^{-1},
 \qquad \HH_h|_K=\sum_i\mathbf h_{K,i}\ell_i\circ\Phi_K^{-1}.
\end{equation}
There is no covariant or contravariant Piola transformation and no built-in
\(H(\mathrm{curl})\) or \(H(\mathrm{div})\) conformity. Curved CAD objects
are represented here by straight tetrahedra and planar boundary triangles.

\subsection{Storage and materials}
\path{MaxwellField} stores six Float64 matrices \path{Ex,Ey,Ez,Hx,Hy,Hz},
each \(N_p\times N_{\rm local}\); \path{MaxwellRHS} mirrors these arrays.
For example \path{U.Hy[3,7]} is the y-component at local interpolation node
3 of local element 7, not a modal coefficient or global node value. DG
coefficients remain distinct even when physical coordinates coincide.
\src{src/solver/data_containers.jl}{MaxwellField; MaxwellRHS; MaxwellElementMaterials}{337--401}

Both scalar homogeneous parameters and element-indexed material vectors
exist. They divide the local electric and magnetic residuals by
\(\varepsilon_K\) and \(\mu_K\), equivalent to inversion of
\(\varepsilon_KM_K\) and \(\mu_KM_K\) for constant coefficients. The
Poisson-bracket flux does not use impedance weighting from the neighboring
material. No anisotropic tensor or spatially varying within-element weighted
mass is represented by this API. Material-ID lookup is merely a means of
constructing elementwise scalar arrays.
\src{src/PhysicalCoverage.jl}{maxwell_element_materials; maxwell_volume_rhs!}{1--79, 332--401}

\section{Exact electric and magnetic residuals}\label{sec:residual}
Define, for three scalar derivative/stiffness matrices,
\begin{equation}\label{eq:curlblock}
 \mathcal C(X,Y,Z)=
 \begin{bmatrix}0&-Z&Y\\Z&0&-X\\-Y&X&0\end{bmatrix},
 \qquad K_K=j_K\mathcal C(S_{K,x},S_{K,y},S_{K,z}).
\end{equation}
The transpose identity is
\(\mathcal C(X^T,Y^T,Z^T)=-\mathcal C(X,Y,Z)^T\), because the cross-component
signs also transpose. The magnetic code has a \emph{positive} addition of
the curl built from \(S^T\), and this is exactly \(-K_K^T\), not \(+K_K^T\).
\src{src/solver/kernels/maxwell.jl}{weak_curl_element!; maxwell_volume_rhs!(PoissonBracketFormulation)}{841--866, 997--1066}

Write \(T_f=I_3\otimes R_f\),
\(W_f^K=I_3\otimes(s_{Kf}W_f)\), and let \(N_f\) act nodewise as
\(\n_K\times\). The complete local coefficient equations are
\begin{align}
 \varepsilon_K(I_3\otimes M_K)\dot e_K
 &=K_Kh_K+\sum_f T_f^TW_f^K N_f(\widehat H-H^-),\label{eq:localE}\\
 \mu_K(I_3\otimes M_K)\dot h_K
 &=-K_K^Te_K-\sum_f T_f^TW_f^K N_f\widehat E.\label{eq:localH}
\end{align}
Here \(\widehat E,\widehat H\) denote vectors of face-node values. Stored flux
buffers already include \(\varepsilon_K^{-1}\) or \(\mu_K^{-1}\); equations
\eqref{eq:localE}--\eqref{eq:localH} have multiplied those factors back to
the left. The lifting subsequently applies \eqref{eq:geomlift} once.

Testing these matrices against coefficient vectors gives the exact weak
description (under the integral/matrix consistency assumptions above):
\begin{align}
 (\varepsilon\dot\E_h,v)_K
 &=(\curl\HH_h,v)_K+
 \langle\n_K\times(\widehat\HH_h-\HH_h^-),v\rangle_{\partial K}\notag\\
 &=(\HH_h,\curl v)_K+\langle\n_K\times\widehat\HH_h,v\rangle_{\partial K},\label{eq:weakE}\\
 (\mu\dot\HH_h,w)_K
 &=-(\E_h,\curl w)_K-\langle\n_K\times\widehat\E_h,w\rangle_{\partial K}\notag\\
 &=-(\E_h,\curl w)_K+\langle\widehat\E_h,\n_K\times w\rangle_{\partial K}.\label{eq:weakH}
\end{align}
The second electric line uses \eqref{eq:curlgreen}. The magnetic surface
term contains the full electric trace, not its jump correction, because
the magnetic volume was already weak. Equivalently, using integration by
parts on \eqref{eq:weakH} gives strong magnetic curl plus
\(-\n\times(\widehat\E-\E^-)\). Mixing this strong correction with the
implemented weak volume would count a boundary term incorrectly.

\subsection{Centered and alternating traces}\label{sec:traces}
Both supported conservative choices can be written with
\(\alpha=1-\theta\):
\begin{equation}\label{eq:hats}
 \widehat\E=\alpha\E^-+\theta\E^+,\qquad
 \widehat\HH=\theta\HH^-+\alpha\HH^+.
\end{equation}
Centered flux sets \(\theta=1/2\). The homogeneous alternating implementation
sets
\begin{equation}\label{eq:theta}
 \theta(\n)=\begin{cases}1,&n_x+n_y+n_z>0,\\
 0,&n_x+n_y+n_z<0,\\1/2,&n_x+n_y+n_z=0.\end{cases}
\end{equation}
Thus faces with exactly zero sum revert to centered traces; the comparison
has no tolerance. This is a fixed spatial orientation selection, not
alternation in time and not selection by MPI rank or element number.
There is no independently adjustable \(C_{12}\) or \(\beta\) parameter.
One can derive \(\widehat\E=\{\E\}+(1/2-\theta)\Delta\E\) and
\(\widehat\HH=\{\HH\}-(1/2-\theta)\Delta\HH\), but this does not introduce
a new implemented user parameter.

The actual buffer values are
\begin{equation}\label{eq:buffers}
 g_E=-\frac{\alpha}{\varepsilon_K}\n\times(\HH^--\HH^+),
 \qquad g_H=-\frac1{\mu_K}\n\times(\alpha\E^-+\theta\E^+).
\end{equation}
\src{src/solver/kernels/maxwell_fluxes.jl}{poisson_bracket_alternating_theta; maxwell_poisson_bracket_surface_flux_values!}{167--180, 251--306}

\subsection{Worked two-sided interface and permutation}
Let \(\n=(1,0,0)\) point from element minus to plus. For alternating flux,
\(\theta^-=1\), so \(\widehat E=E^+\), \(\widehat H=H^-\).
The minus electric face buffer vanishes, while
\(g_H^-=-\mu_-^{-1}(0,-E_z^+,E_y^+)\).
On the plus side \(\n^+=-\n\), \(\theta^+=0\), and its own minus state is
the original plus state. Hence
\(g_E^+=\varepsilon_+^{-1}\n\times(H^+-H^-)\) and
\(g_H^+=\mu_+^{-1}\n\times E^+\). The same physical traces were selected
from both sides, with opposite surface orientation.

If \(P_{ij}=1\) when \(j=\texttt{perm}[i]\), then plus values in minus
ordering are \(P R_+u_+\). The code computes both sides in that ordering,
then scatters the plus buffer using the inverse action \(P^T\), before the
plus lift. In exact arithmetic a matching physical face satisfies
\(s_-W_-=s_+PW_+P^T\). This is the face-mass compatibility needed below.
\src{src/solver/kernels/trace_maps.jl}{match_face_node_permutation; build_dg_trace_maps}{54--170}
\src{src/solver/kernels/maxwell.jl}{interior_face_traces!; unpermute_plus_face_values!}{356--376, 1159--1169}
\src{src/solver/kernels/maxwell_fluxes.jl}{maxwell_interior_surface_face_rhs!(PoissonBracketFormulation)}{882--948}

\subsection{Dispatch discrepancy: material-aware traces}\label{sec:dispatch}
The wrappers \path{maxwell_material_face_flux!} and its allocating counterpart
accept \path{PoissonBracketFormulation} but omit
\path{flux_kind = formulation.flux_kind} when calling the Poisson-bracket
flux routine. Its default is centered. Therefore the element-material path
uses \(\theta=1/2\) even when its formulation object requests alternating
flux. The same wrapper is called by material interior, boundary and
distributed-periodic surface assembly. This affects constant material arrays
as well as heterogeneous arrays. It does not affect the scalar homogeneous
periodic path, which passes the selector explicitly.

The report describes both executed operators. It does not silently insert
the missing keyword. In particular, tests using the material overload under
an alternating label are not evidence that true alternating traces passed.
\src{src/PhysicalCoverage.jl}{maxwell_material_face_flux!; maxwell_material_face_flux}{254--274, 421--437}
\src{src/PhysicalCoverage.jl}{add_material_interior_surface_rhs!; add_distributed_periodic_surface_rhs!}{439--508, 1361--1447}

\section{Boundary conditions and forcing}\label{sec:bc}
All boundary traces below use the outward normal of the computational
material domain, including cavity/obstacle faces. Write
\(v_t=v-\n(\n\cdot v)\) and \(Z=\sqrt{\mu/\varepsilon}\).
\begin{longtable}{p{0.15\linewidth}p{0.75\linewidth}}
\toprule Kind & Executed exterior state or action\\\midrule
PEC & \(E^+=\n(\n\cdot E^-)-E_t^-\), \(H^+=H^-\).\\
PMC & \(E^+=E^-\), \(H^+=\n(\n\cdot H^-)-H_t^-\).\\
Absorbing & \(E^+=\n(\n\cdot E^-)-Z\n\times H^-\),
\(H^+=\n(\n\cdot H^-)+Z^{-1}\n\times E^-\).\\
None & No boundary face residual is added. This is not an outgoing or transparent boundary condition.\\
\bottomrule
\end{longtable}
\src{src/solver/kernels/maxwell.jl}{pec_boundary_plus_trace!; pmc_boundary_plus_trace!; absorbing_boundary_plus_trace!}{411--514}

\subsection{Homogeneous PEC and PMC}
With centered traces, PEC gives \(\n\times\widehat E=0\) and
\(\widehat H=H^-\). Both face buffers in \eqref{eq:buffers} vanish.
PEC is still represented: the magnetic equation retains its weak volume,
equivalent to the strong boundary correction \(+\n\times E^-\).
Setting buffers to zero is not the same as removing PEC from a strong-form
magnetic operator. The homogeneous scalar alternating boundary routine also
returns immediately for PEC, overriding the interior orientation formula.

Centered PMC gives \(\widehat E=E^-\),
\(\n\times\widehat H=0\), so buffers are
\(-\varepsilon^{-1}\n\times H^-\) and
\(-\mu^{-1}\n\times E^-\). The homogeneous scalar alternating path rejects
PMC and absorbing boundaries. By contrast the material path does not apply
this restriction and uses centered traces because of Section~\ref{sec:dispatch}.
These are dispatch-dependent facts, not an assertion of general alternating
PMC support.
\src{src/solver/kernels/maxwell_fluxes.jl}{maxwell_boundary_surface_face_rhs!(PoissonBracketFormulation)}{1199--1274}

\subsection{Absorbing boundary power}
Substitution of the absorbing ghost into centered flux gives
\begin{align}
 \varepsilon g_E&=-\tfrac12\n\times H^- -\tfrac1{2Z}E_t^-,\\
 \mu g_H&=-\tfrac12\n\times E^- -\tfrac Z2H_t^-.
\end{align}
The skew cross terms cancel in the energy pairing, leaving boundary loss
\begin{equation}\label{eq:absorbing}
 \dot{\mathcal H}_h\big|_{\rm absorbing}
 =-\int_{\partial\Omega_{\rm abs}}
 \left(\frac{|E_t^-|^2}{2Z}+\frac Z2|H_t^-|^2\right)dS\le0.
\end{equation}
Integration here is the face mass pairing of the nodal tangential fields.
This is a dissipative termination, not a closed skew bracket.

\subsection{Incident-aware PEC and time dependence}
On a scattered-field PEC boundary, the supplied incident electric field
changes the exterior state by \(-2E_{{\rm inc},t}\). Therefore centered
\(\widehat E_t=-E_{{\rm inc},t}\), \(\widehat H=H^-\), and the magnetic
buffer gains \(\mu^{-1}\n\times E_{\rm inc}\). This weakly imposes
\(\n\times(E_{\rm scat}+E_{\rm inc})=0\); it is not a nodal constraint
on the interior polynomial. Boundary power is
\(\int H^-\cdot(\n\times E_{\rm inc})\,dS\).

Incident data are evaluated at physical face nodes at the mutable time
in \path{MaxwellBoundaryData}; the face mass integrates their degree-p
interpolant, not exact trigonometric boundary data. Scattering drivers use
the material-aware RHS with these data. An alternating label there currently
inherits the centered dispatch discrepancy, rather than proving a distinct
incident alternating implementation.
\src{src/solver/kernels/maxwell.jl}{MaxwellBoundaryData; pec_incident_boundary_plus_trace!}{265--315, 516--609}
\src{src/PhysicalCoverage.jl}{add_material_boundary_surface_rhs!; maxwell_rhs!}{510--640}

The generic ESPRK object stores composition weights, not classical RK
abscissas. Its generic step repeatedly evaluates the full RHS and updates
only the active field partition. It carries no physical-time argument.
The sphere driver's timed wrapper maintains two clocks. Before the first
partition's stage \(i\), it sets
\(\tau_a=\tau_n+\Delta\tau\sum_{j<i}a_j\); before the second partition it
sets \(\tau_b=\tau_n+\Delta\tau\sum_{j<i}b_j\). Each clock advances after
its own nonzero update. For H-first splitting these are respectively the
H and E clocks. Their values need not be monotone if weights are negative.
The wrapper resets boundary time to \(\tau_n+\Delta\tau\) at the end.
No classical \(c_i\) table should be inferred from these weights.
\src{src/solver/kernels/time_integration.jl}{partitioned_symplectic_rk_step!}{535--578}
\src{examples/distributed_metallic_sphere_scattering.jl}{timed_partitioned_symplectic_rk_step!}{1288--1335}

\subsection{Periodic boundaries and PML are different mechanisms}
Periodic boundaries are paired interfaces: no physical exterior state is
manufactured. Plus points are translated to minus coordinates. Cartesian
field components are unchanged by that translation. General rotations,
nonmatching interpolation/mortars and curved periodic mappings are not
implemented by this pairing. Periodic IDs must have kind None in the
physical-boundary registry to avoid adding a second boundary residual.

The nonlinear PML is a separate nodal source addition to the Maxwell RHS,
not a periodic trace or a new face boundary enum. Its coefficients and
nonlinearity invalidate the closed constant-bracket argument for the full
PML dynamics. This report reconstructs the Maxwell spatial part used inside
that path, without assigning Hamiltonian conservation to the PML sources.
\src{src/PhysicalCoverage.jl}{pair_distributed_periodic_descriptors; maxwell_rhs_periodic!}{995--1070, 1449--1560}
\src{src/NonlinearPML.jl}{add_maxwell_nonlinear_pml_source!; maxwell_nonlinear_pml_rhs!}{280--482}

\section{Global assembly, bracket and energy}\label{sec:global}
Let \(m=3N_pN_K\) count vector coefficients on all \emph{owned global}
elements, without duplicate ghosts, and \(u=(e,h)\in\R^{2m}\). With
consistent matching faces, homogeneous PEC/PMC or periodic closure, and
either the true conservative trace pair or the actually executed centered
material pair, the assembled system is
\begin{equation}\label{eq:global}
 \mathcal M\dot u=\mathcal A u,\quad
 \mathcal M=\begin{bmatrix}M_E&0\\0&M_H\end{bmatrix},\quad
 \mathcal A=\begin{bmatrix}0&C\\-C^T&0\end{bmatrix}.
\end{equation}
Here \(M_E=\operatorname{blockdiag}_K(\varepsilon_K I_3\otimes j_KM)\),
similarly for \(M_H\); storage ordering can differ by a fixed permutation.
\(C\) is the complete electric mass-weighted volume-plus-face operator,
not an unweighted differentiation matrix. In the user's \(C_{HE}\) notation,
\(C_{HE}=-C^T\). Boundary forcing adds a vector \(f(\tau)\) on the right;
absorbing boundaries additionally add symmetric negative blocks.

\begin{proposition}[Closed-interface energy identity]\label{prop:energy}
Assume affine geometry, positive elementwise scalar materials, a fixed
conforming mesh with matching face masses/permutations, consistent opposite
normals, and homogeneous PEC/PMC or translational periodic closure.
For the trace pair \eqref{eq:hats} with \(\theta(-\n)=1-\theta(\n)\), the
real-arithmetic operator satisfies \(\mathcal A^T=-\mathcal A\).
\end{proposition}
\begin{proof}
Volume blocks already have the pairing \((K_K,-K_K^T)\), so their energy
production cancels without further integration by parts. On one interior
face, align all values in minus ordering. Set \(\Delta H=H^--H^+\).
The sum of electric face power from both elements is
\begin{equation}
 -\langle\alpha E^-+\theta E^+,\n\times\Delta H\rangle_f
 =-\langle\widehat E,\n\times\Delta H\rangle_f.
\end{equation}
Magnetic face power is
\(-\langle\Delta H,\n\times\widehat E\rangle_f
=\langle\widehat E,\n\times\Delta H\rangle_f\), and cancels it.
The equal physical face pairing on both sides is essential. The same proof
applies across translated periodic interfaces. PEC face buffers vanish.
PMC cross terms cancel pointwise. Since the operator is linear and its
quadratic energy production vanishes for every coefficient vector, its
symmetric part vanishes; the already explicit complementary field coupling
then yields the block identity in \eqref{eq:global}.
\end{proof}
This argument separates algebraic energy cancellation from the stronger
geometric integration-by-parts identities needed for consistency with the
continuous differential operator. Floating-point material division and
mass solves perturb both statements slightly.

\subsection{Gradients and the finite-dimensional Poisson tensor}
The implemented mass Hamiltonian is
\begin{equation}\label{eq:hamiltonian}
 \mathcal H_h(u)=\tfrac12u^T\mathcal M u,
 \qquad \nabla_u\mathcal H_h=\mathcal M u.
\end{equation}
If a functional has coefficient gradient \(\nabla_e F\), the coefficients
of its unweighted discrete \(L^2\) gradient are \(M_0^{-1}\nabla_e F\),
and those of its material-weighted gradient are \(M_E^{-1}\nabla_e F\).
The corresponding magnetic conversion uses \(M_H^{-1}\). This follows
by equating \(dF=\nabla_eF^T\delta e\) to the chosen mass inner product.

The correct coefficient Poisson tensor is therefore
\begin{equation}\label{eq:J}
 J_h=\mathcal M^{-1}\mathcal A\mathcal M^{-1}
 =\begin{bmatrix}
 0&M_E^{-1}CM_H^{-1}\\
 -M_H^{-1}C^TM_E^{-1}&0
 \end{bmatrix},\quad
 \dot u=J_h\nabla_u\mathcal H_h.
\end{equation}
It is not \(\mathcal M^{-1}\mathcal A\) itself. Define
\(\{F,G\}_h=\nabla_u F^TJ_h\nabla_u G\).
For \(L=\mathcal M^{-1}\mathcal A\), the equivalent weighted adjoint
identity is \(\mathcal ML+L^T\mathcal M=0\).
\src{src/solver/kernels/maxwell.jl}{maxwell_energy; maxwell_energy_rate}{1868--1921, 2011--2037}
\src{src/PhysicalCoverage.jl}{maxwell_energy; distributed_maxwell_energy}{712--799}

\begin{proposition}[Jacobi under fixed closed linear assumptions]
Under Proposition~\ref{prop:energy}, with coefficients, geometry and face
selection fixed independently of \(u\), \eqref{eq:J} defines a constant
skew matrix. Its finite-dimensional bracket satisfies the Jacobi identity.
\end{proposition}
\begin{proof}
The coordinate Jacobi condition is
\(\sum_\ell(J_{i\ell}\partial_\ell J_{jk}
+J_{j\ell}\partial_\ell J_{ki}+J_{k\ell}\partial_\ell J_{ij})=0\).
Every derivative is zero because the matrix is constant. Skew-symmetry,
bilinearity and the Leibniz rule follow directly from the coefficient bracket.
\end{proof}
This is a conditional reconstruction of the fixed linear discretization,
not a proof that every feature bearing the name Poisson bracket, including
nonlinear PML or open boundaries, has this structure. The floating-point
stored/evaluated operator need only approximate a skew matrix. A numerical
skew check by itself is not a proof of a state-dependent Jacobi identity.

\subsection{What time stepping can and cannot inherit}
The conservative ESPRK update composes electric and magnetic shears.
For the autonomous block system, each subflow changes only its own field
using the opposite field, consistent with a split Hamiltonian. Exact
preservation of quadratic energy by the \emph{composition} does not follow
from semidiscrete skew-adjointness. Symplectic/Poisson structure, truncation
error, long-time bounded energy oscillation and exact energy conservation
are distinct assertions. Conversely, an exact linear null invariant of each
partition survives exact-arithmetic linear updates; Float64 updates can
accumulate small signed residuals. Dissipative source additions do not acquire
symplecticity merely by using the same stepping routine.

\section{Distributed-memory realization}\label{sec:mpi}
\subsection{Mesh and ownership}
The legacy VTK loader extracts four-vertex tetrahedra. The periodic driver
reconstructs boundary triangles by counting sorted triples of tetrahedron
vertices, infers box bounds, and labels xmin/xmax/ymin/ymax/zmin/zmax as
1/2/3/4/5/6. It does not rely on a special PBC tag in the VTK file. Mesh
partition files contain zero-based rank assignments, distinct from Julia's
one-based local indices. Root distribution and cached per-rank mesh loading
are both implemented.
\src{src/MetisIO.jl}{read_mesh_file_tet_vtk}{27--163}
\src{src/MaxwellDriverSupport.jl}{load_periodic_box_mesh; periodic_box_specs}{250--310}
\src{src/DistributedIO.jl}{prepare_distributed_mesh_partition; build_distributed_dg_from_partition}{218--288, 365--385}

An element belongs to one rank. Each rank also stores a face-adjacent halo.
Local stored elements are sorted by global element ID, so owned elements
must be accessed through \path{partition.owned}, not assumed to be the first
columns. Vertex order within a copied tetrahedron is preserved. Local topology
is rebuilt on owned and halo elements; ordinary partition interfaces then
appear as local interior faces between an owned and ghost element. The
distributed and solver face-numbering tables are reconciled by sorted global
vertex keys, not by assuming equal face IDs.
\src{src/DistributedMesh3D.jl}{build_ghost_global_elements_face_halo; build_distributed_mesh}{291--319, 762--823}
\src{src/DistributedDG.jl}{local_mesh_from_distributed_mesh; find_solver_face; build_distributed_maxwell_exchange}{81--232}

\subsection{An ordinary remote face, end to end}
Suppose rank A owns element \(K\), and its neighbor \(L\) is owned by B.
The exchange table identifies their common face by the sorted global vertex
triple. Neighbor face lists are sorted by this key. Rank B packs, face by
face, all six Cartesian field components, each in its native face-node order:
\begin{equation}
 [E_x(1{:}N_{fp}),E_y(1{:}N_{fp}),E_z(1{:}N_{fp}),
 H_x(1{:}N_{fp}),H_y(1{:}N_{fp}),H_z(1{:}N_{fp})].
\end{equation}
Thus each directed interface transmits \(6N_{fp}\) Float64 values, not full
element coefficients and not just tangential components. Rank A receives
them into the same face-node ordering of its copied halo element \(L\).
The ordinary local trace map then applies \(P\) to align \(L\)'s values with
\(K\)'s face. This separation explains why unpacking itself has no extra
permutation: the copied element has the original vertex ordering, and the
solver trace map handles cross-element orientation.

All receives and sends are nonblocking MPI calls with tag 24017, followed
by \path{MPI.Waitall} before residual evaluation. There is no overlap of
this exchange with volume computation in this path. The face kernel adds
both local-element contributions, possibly including one ghost contribution.
After assembly, all ghost RHS columns are zeroed. Only the owned contribution
represents a global degree of freedom; no residual scatter-add to the other
rank is required because that rank assembles its own owned side.
\src{src/DistributedDG.jl}{pack_maxwell_interface!; unpack_maxwell_interface!; exchange_maxwell_ghost_traces!; zero_ghost_maxwell_rhs!}{620--731}

\subsection{Periodic remote faces}
Periodic descriptors are gathered to root and paired by translated
centroids; node matching then constructs a bijective permutation, with
area compatibility checked. Default tolerances are \(10^{-9}\) for
centroids/nodes and relative \(10^{-10}\) for area. Pairing is broadcast.
These routines store normals but do not explicitly reject nonopposite
normals in the pairing function: the mathematical cancellation assumes
valid oriented translated faces, and the independent checks must test it.

Each rank stores directed descriptors only for its owned periodic faces.
Same-rank partners are gathered directly. Remote partners communicate the
same six face-component arrays, ordered by pair ID, with tag 24018 and
offset \((i-1)6N_{fp}+1\). The received values are permuted into local
face order before evaluating \eqref{eq:buffers}; only the local owned side
is lifted. The inverse permutation is stored for the opposite side.
If using material arrays, partner material values are included in setup
metadata rather than exchanged every RHS.
\src{src/PhysicalCoverage.jl}{build_distributed_periodic_maxwell_exchange; exchange_distributed_periodic_traces!; distributed_periodic_plus_trace!}{1083--1287, 1323--1359}

Every RHS evaluation exchanges both E and H traces, even when an ESPRK
substep will use only one field's RHS. The distributed periodic RHS performs
ordinary exchange, periodic exchange, volume, ordinary faces, periodic
faces, physical boundaries and ghost zeroing in that order. Global diagnostic
sums include only owned elements and use MPI reductions. Equivalent global
algebra across partitions is not a promise of bitwise reproducibility:
summation order, face ordering and reduction trees can change rounding.
\src{src/PhysicalCoverage.jl}{maxwell_rhs_periodic!}{1449--1560}

\section{Divergence, charge, skeleton terms and null invariants}\label{sec:charge}
Define the broken magnetic divergence coefficient map on each element as
\begin{equation}
 D_{b,K}h_K=\mu_K[D_{K,x}\;D_{K,y}\;D_{K,z}]h_K.
\end{equation}
Its squared L2 norm is \(\sum_K(D_{b,K}h_K)^Tj_KM(D_{b,K}h_K)\).
Production quadrature diagnostics interpolate the same strong derivatives
to integration points, sum squared values, and take a square root. Their
reported infinity norm is a sampled maximum, not a certified continuous
supremum. Neither derivative includes a numerical normal-flux lift.

The signed magnetic-charge functional is exactly
\begin{equation}\label{eq:charge}
 Q_H=\sum_K\mu_Kj_K\mathbf1^TM
 [D_{K,x}\;D_{K,y}\;D_{K,z}]h_K=c_H^Th=q_H^TM_Hh,
 \qquad q_H=M_H^{-1}c_H.
\end{equation}
This is what \path{distributed_strong_magnetic_charge} computes, followed
by an owned-element MPI sum. Evaluating the same routine on the magnetic
RHS gives
\begin{equation}\label{eq:qnull}
 \dot Q_H=c_H^TM_H^{-1}C_{HE}e=q_H^TC_{HE}e.
\end{equation}
The global null condition is \(q_H^TC_{HE}=0\); the much stronger local
condition is \(D_bM_H^{-1}C_{HE}=0\). A small value of the signed expression
for a few electric states does not prove either full matrix identity.
\src{src/MaxwellDiagnostics.jl}{local_strong_magnetic_charge_from_components; distributed_strong_magnetic_charge_derivative}{217--297}

By elementwise divergence theorem, \eqref{eq:charge} can alternatively be
written as the sum of interior/periodic \emph{normal jumps} plus physical
boundary normal fields. This is a representation of the same signed
functional, not an additional correction that should be added to it. A
different DG divergence defined through a single-valued numerical normal
trace is not specified by the tangential Maxwell flux routines. It must
not be invented and then attributed to the implementation.

The RHS breakdown separately applies this signed functional to the volume,
ordinary-face, periodic-face and physical-boundary RHS pieces. Its ``face
contributions'' are contributions to \(\dot Q_H\), not an independently
defined skeleton charge that repairs \(Q_H\). It currently uses the material
overload, so the flux-dispatch caveat also applies to alternating-labeled
breakdown tests.
\src{src/MaxwellDiagnostics.jl}{distributed_magnetic_charge_rhs_breakdown}{299--415}

\subsection{What would qualify as a Casimir?}
For a coefficient-linear functional \(C_0(u)=c^Tu\), being a Casimir of
\eqref{eq:J} requires \(c^TJ_h=0\), independently of the Hamiltonian.
Conservation only along a particular trajectory is weaker. Conservation for
\emph{all} states under the quadratic Hamiltonian here, however, means
\(c^TJ_h\mathcal M u=0\) for every \(u\); because \(\mathcal M\) is
invertible, it implies \(c^TJ_h=0\) under the closed constant-skew assumptions.
The distinction is therefore between a verified full identity and a few
observed trajectories, as well as between closed and forced/dissipative paths.
This report calls the monitored quantity a \emph{candidate null invariant}
unless that identity is established for the stated mesh and operator.

Curl-potential initialization imposes \(D_bh(0)\simeq0\) on affine elements
but does not change \(D_bM_H^{-1}C_{HE}\). The long-time run documented in
\path{docs/numerics/magnetic_charge_longtime_analysis.md} exhibits precisely
this distinction: local divergence grows to order one while signed charge
drifts only slightly. That report's checkpoint/replay numbers are separate
evidence, not rerun measurements in this operator audit.
\src{src/MaxwellInitialization.jl}{initialize_magnetic_from_potential!; magnetic_initialization_diagnostics}{41--169}

Momentum, angular momentum and chirality use further quadratic/derivative
functionals, not just the Hamiltonian matrix. Their preservation requires
their own algebraic identities. In the quadrature diagnostics momentum is
the integral of \(\varepsilon\mu\E\times\HH\), angular momentum includes
the position cross product, and chirality uses the broken curls in
\(\tfrac12(\varepsilon\E\cdot\curl\E+\mu\HH\cdot\curl\HH)\).
Antisymmetry of \(J_h\) alone does not preserve these functionals. Position
weighting also requires care with periodic-cell boundary fluxes.

\section{A fully explicit degree-one example}\label{sec:p1}
Use the equispaced ordering \(v_4,v_3,v_2,v_1\). The nodal row basis is
\begin{equation}
 \ell=\tfrac12(t+1,\ s+1,\ r+1,\ -1-r-s-t).
\end{equation}
Its reference matrices, derived by simplex moments, are
\begin{equation}\label{eq:p1mass}
 M=\frac1{15}\begin{bmatrix}2&1&1&1\\1&2&1&1\\1&1&2&1\\1&1&1&2\end{bmatrix},
 \quad M^{-1}=15I-3\mathbf1\mathbf1^T,
\end{equation}
\begin{equation}
 D_r=\mathbf1[0,0,1/2,-1/2],\quad
 D_s=\mathbf1[0,1/2,0,-1/2],\quad
 D_t=\mathbf1[1/2,0,0,-1/2].
\end{equation}
Each row of \(S_r=MD_r\) is \([0,0,1/6,-1/6]\), similarly for the other
directions. Notice the factor \(1/3\) from the mass row sum.

For face 4, the ordered volume node IDs are \((3,2,4)\), corresponding to
triangle nodes \((a,b)=(-1,1),(1,-1),(-1,-1)\). Its matrices are
\begin{equation}
 R_4=\begin{bmatrix}0&0&1&0\\0&1&0&0\\0&0&0&1\end{bmatrix},\quad
 W_4=\frac16\begin{bmatrix}2&1&1\\1&2&1\\1&1&2\end{bmatrix},\quad
 L_4=\begin{bmatrix}-2&-2&-2\\1/2&3&1/2\\3&1/2&1/2\\1/2&1/2&3\end{bmatrix}.
\end{equation}
For example \(L_4\mathbf1=(-6,4,4,4)^T\); a constant face forcing affects
all four volume coefficients, including the off-face node.

Take physical vertices \((0,0,0),(2,0,0),(0,1,0),(0,0,1)\). Then
\(J=\operatorname{diag}(1,1/2,1/2)\), \(j=1/4\),
\(D_x=D_r,D_y=2D_s,D_z=2D_t\), and \(M_K=M/4\).
Face 4 has physical area 1, so \(s_{K4}=1/2\) and its physical lift is
\(2L_4\). All these numbers are hand-derived for this mapping.

For \(H=(0,0,y)\), the nodal \(H_z\) values are \((0,1,0,0)^T\).
The electric x-volume residual is
\(\varepsilon^{-1}M^{-1}(2S_s)H_z
=\varepsilon^{-1}\mathbf1\), agreeing with \(\curl H=(1,0,0)\).
For \(E=(0,0,y)\), the weak magnetic x-volume residual is instead
\begin{equation}
 \mu^{-1}M^{-1}(2S_s^T)E_z
 =\mu^{-1}(0,5,0,-5)^T.
\end{equation}
It must \emph{not} be mistaken for the complete negative strong curl.
If the numerical electric trace equals the interior trace on all faces,
the missing surface terms change this vector to
\(-\mu^{-1}\mathbf1\), by \eqref{eq:sbp}. This explicit example exposes
why the magnetic volume and surface terms must be audited together.

\section{Independent verification}\label{sec:verification}
The accompanying \path{verify_operator.jl} is a diagnostic, not a production
replacement. It constructs the matrices in \eqref{eq:localE}--\eqref{eq:localH}
independently of the production flux and lift routines, then applies the
actual matrix-free RHS to every coefficient basis vector on six- and
twelve-tetrahedron unit cubes. The assembled comparison is \(\mathcal ML_{\rm actual}\) versus
the independently reconstructed \(\mathcal A\). This checks every column
for these small problems, rather than a single random vector.

Reference checks cover both nodal families at degrees 1--4, constants,
coordinates, positive mass eigenvalues, integration by parts, face moments,
and skew affine maps with both determinant signs. Diagnostic cubature checks
cover all unit-simplex monomials up to orders 2, 4 and 8. Small global checks
cover degrees 1 and 2, periodic centered/alternating, PEC centered/alternating,
PMC centered, and absorbing centered on six tetrahedra. The twelve-element
periodic cases add interior faces on which the alternating selector is
noncentral; a degree-one heterogeneous PEC case checks material-weighted
mass inversion. Distributed checks use twelve tetrahedra, degrees 1--3,
and compare all owned
RHS coefficients against serial execution with poisoned ghost coefficients,
using both nodal families and true homogeneous flux selections. They also
measure the material-wrapper discrepancy separately. The serial comparisons
are globally reduced maxima; material-dispatch difference rows in the CSV
are measurements on rank zero's local arrays, not global maxima.

Reference residual tolerances are
\(2000\epsilon_{\rm mach}\max(1,\kappa_2(V)^2)\), increased by five for the
skew physical-map checks. Matrix reconstruction and skew tests use
\(3\times10^{-12}\) after scaling by \(\max(1,\|\mathcal A\|_\infty)\).
The MPI owned-entry maximum difference has tolerance \(2\times10^{-10}\).
These are deliberately conservative small-test tolerances, not measured
errors. The charge row residual is normalized by
\(\max(1,\|c_H^T\|_1\|L_{HE}\|_\infty)\); local divergence uses the
unscaled induced infinity norm \(\|D_bL_{HE}\|_\infty\).
Measured values and execution status are recorded in the CSV files and audit.
A CSV tolerance of NaN denotes a measurement, not a pass/fail criterion;
the corresponding \texttt{pass=true} is vacuous and must not be read as
an assertion that, for example, the broken divergence operator is null.
The pre-existing reference-node test file also passed all 761 assertions.
It emitted a modal-orthonormality warning of \(1.38\times10^{-9}\) during
its higher-degree construction; the saved log retains that warning.
\subsection{Executed results}
All three independent invocations (1, 2 and 4 MPI ranks) completed with
exit status zero. Each wrote 173 measurements, of which 110 have explicit
tolerance checks; all 110 passed. Positive-definiteness is additionally
asserted in the script. These counts exclude the separate pre-existing tests.
The tables below summarize the saved CSV values, not predicted tolerances.
Float64 machine epsilon is \(2.220446049250313\times10^{-16}\).

\begin{center}
\begin{tabular}{lrr}
\toprule
Reference check & Equispaced & Warp-blend\\\midrule
\(p=1\) IBP residual & \(5.27\times10^{-16}\) & \(5.14\times10^{-16}\)\\
\(p=2\) IBP residual & \(3.23\times10^{-15}\) & \(3.55\times10^{-15}\)\\
\(p=3\) IBP residual & \(5.36\times10^{-13}\) & \(4.69\times10^{-13}\)\\
\(p=4\) IBP residual & \(6.12\times10^{-11}\) & \(5.06\times10^{-11}\)\\
\(p=4\) face moment cross-check & \(1.44\times10^{-11}\) & \(1.03\times10^{-11}\)\\
\bottomrule
\end{tabular}
\end{center}
The IBP entries are the maximum induced infinity norm over the three
reference directions. All mass matrices tested were symmetric with positive
eigenvalues. Constant and coordinate differentiation errors were at most
\(4.28\times10^{-15}\) and \(1.08\times10^{-14}\), respectively; the
nonorthogonal affine coordinate checks reached \(2.29\times10^{-14}\).
Cubature moment errors were at most \(2.23\times10^{-16}\).
The degree-four IBP residual is clearly larger than a few machine epsilons.
Analytic real-arithmetic moments do not imply machine-precision Float64
matrix identities after cancellation-prone Gram construction. The observed
degree dependence is consistent with that conditioning issue, but these
checks do not isolate its contribution from every other construction step.

\begin{center}
\begin{tabular}{p{0.64\linewidth}r}
\toprule
Complete-operator or MPI check & Maximum residual\\\midrule
Scaled independent reconstruction, all tested cases & \(2.68\times10^{-16}\)\\
Scaled weighted skew identity, closed cases & \(8.44\times10^{-16}\)\\
Heterogeneous PEC reconstruction, 12 tets, \(p=1\) & \(1.39\times10^{-16}\)\\
Absolute energy production, closed deterministic test states & \(4.72\times10^{-16}\)\\
MPI owned-RHS difference, 1 rank versus serial & \(0\)\\
MPI owned-RHS difference, 2 ranks versus serial & \(2.84\times10^{-14}\)\\
MPI owned-RHS difference, 4 ranks versus serial & \(2.84\times10^{-14}\)\\
\bottomrule
\end{tabular}
\end{center}
The reconstruction/skew checks assemble every column at degrees 1 and 2.
The MPI checks instead use a deterministic full coefficient vector at
degrees 1--3 with both nodal families and both scalar flux choices. They
exercise ordinary and periodic halo exchange, including poisoned ghost
input, but are not exhaustive distributed column-by-column proofs.
The absorbing symmetric-part eigenvalue test is nonpositive to its
\(3\times10^{-12}\) tolerance; deterministic energy productions are
\(-5.37\) and \(-5.87\) for degrees 1 and 2.

\begin{center}
\begin{tabular}{llrr}
\toprule
Mesh & Case & Scaled charge row residual & \(\|D_bL_{HE}\|_\infty\)\\\midrule
6 tets & \(p=1\), periodic centered & \(6.25\times10^{-18}\) & 50\\
6 tets & \(p=1\), periodic alternating & \(7.71\times10^{-18}\) & 70\\
6 tets & \(p=2\), periodic centered & \(2.51\times10^{-17}\) & 564\\
6 tets & \(p=2\), periodic alternating & \(2.69\times10^{-17}\) & 660\\
12 tets & \(p=2\), periodic centered & \(1.26\times10^{-17}\) & 952\\
12 tets & \(p=2\), periodic alternating & \(1.03\times10^{-17}\) & 1384\\
6 tets & \(p=1\), PMC centered & \(4.08\times10^{-3}\) & 50\\
\bottomrule
\end{tabular}
\end{center}
The global periodic charge identity is consistent with roundoff on these
meshes, while the complete broken divergence-of-curl map is plainly nonzero.
These are compatible findings, not contradictory diagnoses. The PMC row
also demonstrates that this particular charge diagnostic is not a universal
null invariant across boundary closures, even though centered PMC conserves
the Hamiltonian. No charge threshold was enforced in the diagnostic script;
these are measured operator rows, interpreted with their specified scaling.

Every material-aware alternating-versus-centered difference was exactly
zero in the tested arrays, while the difference from the true scalar
alternating operator was nonzero. For example the single-rank degree-three
equispaced case differed by \(1.34\times10^2\) in maximum RHS entry.
This independently corroborates the missing flux-selector dispatch, rather
than attributing that difference to MPI or roundoff.


\section{Discrepancies and limitations}\label{sec:limitations}
\begin{enumerate}
\item Material-aware alternating requests execute centered traces because
the wrapper omits the flux selector. Production code was intentionally left
unchanged. Earlier alternating-labeled material tests cannot establish true
alternating behavior.
\item The alternative face-quadrature helper attempts to broadcast a row
matrix into a vector slice in \path{eval_nodal_basis_at_points}. The audit
records the resulting \path{DimensionMismatch}. Active face mass construction
uses triangular Gram matrices instead and does not call that helper.
\item The older \path{docs/math/maxwell_poisson_bracket_formulation.tex},
line 134, says alternating is not an exposed production option. The current
scalar path does expose it. This report supersedes that particular statement,
without rewriting unrelated documentation.
\item No within-element anisotropic material, curved metric, nonmatching
periodic mortar or Piola-conforming field representation is reconstructed
because none is implemented on this path.
\item Pairing verifies translated nodes/area but not a separate opposite-normal
condition. Valid mesh geometry remains an assumption; passing small MPI tests
does not certify arbitrary supplied meshes.
\item The monitored signed charge and local broken divergence are different
functionals. No universal local divergence-preservation claim follows from
energy or signed-charge cancellation. General-mesh charge-null proofs remain
separate from the finite meshes tested here.
\item Closed, constant-coefficient-matrix Jacobi follows from a derivation;
PML, absorbing termination and externally forced evolution are not covered
by that closed Poisson claim. Fully discrete energy conservation is not
claimed for ESPRK.
\end{enumerate}

\appendix
\section{Mathematical-object-to-code map}
All paths below are relative to the repository root. Earlier source blocks
give function names and exact ranges; this table is a compact navigation aid.
\small
\begin{longtable}{>{\raggedright\arraybackslash}p{0.27\linewidth}p{0.65\linewidth}}
\toprule Object / operation & Executable location\\\midrule\endhead
Reference vertices, field layouts & \path{src/solver/data_containers.jl}, 88--164, 337--401.\\
Gram, basis, mass, stiffness & \path{src/solver/kernels/reference_tet.jl}, 426--646, 769--848.\\
Face restriction, face mass, lift & \path{src/solver/kernels/face_operators.jl}, 243--268, 734--972.\\
Affine metrics, physical derivatives & \path{src/solver/kernels/metrics.jl}, 7--98; \path{src/solver/kernels/physical_operators.jl}, 6--80.\\
Normal orientation and areas & \path{src/solver/kernels/geometry.jl}, 113--148.\\
Volume curl and adjoint curl & \path{src/solver/kernels/maxwell.jl}, 841--866, 997--1066.\\
Conservative trace pair & \path{src/solver/kernels/maxwell_fluxes.jl}, 102--180, 251--320.\\
Two-sided assembly & \path{src/solver/kernels/maxwell_fluxes.jl}, 882--948, 1042--1108.\\
Boundary ghosts and incident forcing & \path{src/solver/kernels/maxwell.jl}, 411--609.\\
Material dispatch & \path{src/PhysicalCoverage.jl}, 254--274, 421--640.\\
Ordinary MPI halo exchange & \path{src/DistributedDG.jl}, 181--232, 620--731.\\
Periodic pairing and assembly & \path{src/PhysicalCoverage.jl}, 995--1560.\\
Signed charge and split RHS diagnostics & \path{src/MaxwellDiagnostics.jl}, 217--415.\\
Diagnostic cubature and invariants & \path{src/MaxwellDiagnostics.jl}, 417--780; \path{src/solver/kernels/JaskowiecSukumar.jl}.\\
ESPRK partition updates & \path{src/solver/kernels/time_integration.jl}, 291--369, 535--578.\\
Timed incident data & \path{examples/distributed_metallic_sphere_scattering.jl}, 1288--1335.\\
\bottomrule
\end{longtable}
\normalsize

\clearpage
\section{Reproducibility and interpretation of the audit}
The accompanying files are \path{implementation_audit.md},
\path{verify_operator.jl}, \path{checks_rank1.csv}, \path{checks_rank2.csv},
\path{checks_rank4.csv}, and \path{provenance.txt}. The hash inventory covers
the source files on which the derivation depends; local modifications are
part of the snapshot. No repository source or test was repaired for this
report. The diagnostic's own development corrections do not modify the
operator under audit.

The executed environment was Julia 1.10.11, MPI.jl 0.20.23 with MPICH 4.3.1
(ch3/nemesis), and ILP64 OpenBLAS. Julia and BLAS each used one thread per
MPI process. The provenance can be refreshed, without modifying solver
files, using the included \path{capture_provenance.py} script; doing so
creates a new snapshot rather than reproducing the old one.

From the repository root, with a matching MPI launcher:
\begin{lstlisting}
export OPENBLAS_NUM_THREADS=1 JULIA_NUM_THREADS=1 OMP_NUM_THREADS=1
julia --project=. --compiled-modules=no \
  docs/math/poisson_bracket_implementation/verify_operator.jl /tmp/pb_rank1.csv
julia --project=. --compiled-modules=no -e '
using MPI
run(`$(MPI.mpiexec()) -n 2 $(Base.julia_cmd()) --project=. --compiled-modules=no docs/math/poisson_bracket_implementation/verify_operator.jl /tmp/pb_rank2.csv`)
'
# Repeat with -n 4 for the four-rank comparison.
\end{lstlisting}
Compile from the document directory with
\begin{lstlisting}
pdflatex -interaction=nonstopmode -halt-on-error poisson_bracket_implementation.tex
pdflatex -interaction=nonstopmode -halt-on-error poisson_bracket_implementation.tex
\end{lstlisting}
These checks audit a fixed linear spatial operator. They are not a full
convergence study, a PML validation, a high-degree conditioning survey or a
proof for every admissible tetrahedral mesh. Numerical verification supports
the source derivation only within its explicitly recorded scope.
\end{document}
